\documentclass[../../main.tex]{subfiles}
\begin{document}
\begin{oframed}
  \subfile{problem}
\end{oframed}

{\bf [解答]}
$c,d$を解に持つ方程式は
\begin{align}
  x^2 - (c+d)x + cd = 0
\end{align}
だから，$c+d,cd$を求めれば良い．
$
\begin{pmatrix} a \\ b
\end{pmatrix}$と$
\begin{pmatrix} c \\ d
\end{pmatrix}$が直交するから
\begin{align}
  ac + bd = 0 \label{eq:1}
\end{align}
であり，また$\vec{v}$は大きさが$1$だから
\begin{align}
  c^2 + d^2 = 1 \label{eq:2}
\end{align}
である．また，$a,b$は$x^2+x-1=0$の解だから解と係数の関係より
\begin{align}
  a+b = -1 , ab = -1 \label{eq:3}
\end{align}
である．よって$ab \neq 0$ だから\cref{eq:1}より，
\begin{align}
  d = -\dfrac{a}{b}c
\end{align}
となり，\cref{eq:2}に代入して
\begin{align}
  \left( 1 + \dfrac{a^2}{b^2} \right) c^2 = 1
  \therefore
  c = \pm \dfrac{b}{\sqrt{a^2+b^2}} = \pm \dfrac{b}{\sqrt{(a+b)^2-2ab}} = \pm\dfrac{b}{\sqrt{3}} \quad (\because\text{\cref{eq:3}})
\end{align}
を得る．\cref{eq:1}に代入して $d = \mp \dfrac{a}{\sqrt{3}}$（以下複号同順）となるから
\begin{align}
  (c+d, cd) = \left( \pm \dfrac{1}{\sqrt{3}}(b-a), \dfrac{1}{3} \right)
\end{align}
である．ここで$x^2+x-1=0$を解いて$a,b=\dfrac{-1 \pm \sqrt{5}}{2}$だから
$b-a = \pm \sqrt{5}$ となるので代入して
\begin{align}
  (c+d, cd) = \left( \pm \dfrac{1}{3}\sqrt{15}, \dfrac{1}{3} \right)
\end{align}
を得る．よって$c, d$ を 2解に持つ方程式は
\begin{align}
  x^2 \pm \dfrac{1}{3}\sqrt{15}x + \dfrac{1}{3} = 0
\end{align}
である．

\newpage

\end{document}
